\chapter{高级有限群论专题}
\label{app:advanced-finite-group}

本附录收束\chapref{chap:homomorphism-quotient} 正文中移出的纯有限群论技术细节：特征子群与 $N/C$ 定理、Frattini 论证、对称群的自同构、Jordan--H\"older 定理与可解群的完整证明、幂零群、自由群与 Tietze 变换、$A_5$ 单性的完整枚举。物理读者掌握主线（共轭类、商群、群作用、轨道--稳定子、Burnside 计数、Sylow 三结论、直积与半直积）即可进入表示论与点群；这里的内容为需要查考代数细节时准备，不要求按序阅读。

本附录沿用全书符号约定：群元换位子为
\[
[x,y]=xyx^{-1}y^{-1},
\]
导群 $G'=[G,G]$ 为全部换位子生成的子群；子群 $H$ 的正规化子为 $N_G(H)=\{g\in G\mid gHg^{-1}=H\}$，中心化子为 $C_G(H)=\{g\in G\mid gh=hg,\ \forall h\in H\}$。下文多次提到 Sylow $p$-子群：若 $|G|=p^am$ 且 $p\nmid m$，$G$ 的 $p^a$ 阶子群称为 Sylow $p$-子群，其存在性、共轭性与个数约束（Sylow 定理）见\cref{thm:sylow}。

\section{特征子群、正规化子与自同构}
\label{sec:char-normalizer}

\begin{definition}[特征子群]
子群 $H\leqslant G$ 称为\textbf{特征子群}（记 $H\operatorname{char}G$），若对所有 $\varphi\in\operatorname{Aut}(G)$ 都有 $\varphi(H)=H$。
\end{definition}

特征子群比正规子群更强：正规子群只对内自同构不变，特征子群对所有自同构不变。

\begin{proposition}[特征子群的性质]
\label{prop:characteristic-subgroup}
\begin{enumerate}
\item 若 $H\operatorname{char}G$，则 $H\mathrel{\trianglelefteq} G$。
\item 若 $H\operatorname{char}K$ 且 $K\operatorname{char}G$，则 $H\operatorname{char}G$（传递性）。
\item 若 $H\operatorname{char}K$ 且 $K\mathrel{\trianglelefteq} G$，则 $H\mathrel{\trianglelefteq} G$。
\item $Z(G)\operatorname{char}G$，$G'\operatorname{char}G$（导群是特征子群）。
\end{enumerate}
\end{proposition}

\begin{proof}
(1) 内自同构是自同构的特例，故 $\varphi(H)=H$ 对所有 $\varphi\in\operatorname{Aut}(G)$ 成立蕴含 $c_g(H)=H$ 对所有 $g\in G$ 成立，即 $gHg^{-1}=H$，故 $H\mathrel{\trianglelefteq} G$。

(2) 对任意 $\varphi\in\operatorname{Aut}(G)$，$\varphi|_K\in\operatorname{Aut}(K)$（因 $K\operatorname{char}G$ 保证 $\varphi(K)=K$），故 $\varphi(H)=\varphi|_K(H)=H$（因 $H\operatorname{char}K$）。

(3) 对任意 $g\in G$，$c_g\in\operatorname{Inn}(G)\subseteq\operatorname{Aut}(G)$，$c_g(K)=K$（因 $K\mathrel{\trianglelefteq} G$），故 $c_g|_K\in\operatorname{Aut}(K)$，$c_g(H)=c_g|_K(H)=H$（因 $H\operatorname{char}K$）。

(4) 对任意 $\varphi\in\operatorname{Aut}(G)$ 和 $z\in Z(G)$：$\varphi(z)x\varphi(z)^{-1}=\varphi(z\varphi^{-1}(x)z^{-1})=\varphi(\varphi^{-1}(x))=x$，故 $\varphi(z)\in Z(G)$，即 $\varphi(Z(G))\subseteq Z(G)$；同理论 $\varphi^{-1}$ 得反向包含。导群 $G'=[G,G]$ 由换位子生成，$\varphi([x,y])=[\varphi(x),\varphi(y)]$，故 $\varphi(G')=G'$。
\end{proof}

\begin{theorem}[$N/C$ 定理]
\label{thm:nc-theorem}
设 $H\leqslant G$，则 $N_G(H)/C_G(H)$ 同构于 $\operatorname{Aut}(H)$ 的一个子群：
\[
N_G(H)/C_G(H)\hookrightarrow\operatorname{Aut}(H).
\]
特别地，$|N_G(H)/C_G(H)|\bigm||\operatorname{Aut}(H)|$。
\end{theorem}

\begin{proof}
定义映射 $\Phi:N_G(H)\to\operatorname{Aut}(H)$，$\Phi(g)=c_g|_H$（$h\mapsto ghg^{-1}$）。对 $g\in N_G(H)$，$c_g(H)=gHg^{-1}=H$，故 $c_g|_H$ 是 $H$ 到自身的同构，即 $\Phi(g)\in\operatorname{Aut}(H)$。

同态性。$\Phi(gh)(x)=ghxh^{-1}g^{-1}=g(hxh^{-1})g^{-1}=\Phi(g)\circ\Phi(h)(x)$。

核。$\Phi(g)=\id(H)$ 当且仅当 $ghg^{-1}=h$ 对所有 $h\in H$，即 $g\in C_G(H)$。故 $\operatorname{Ker}\Phi=C_G(H)$。

由第一同构定理，$N_G(H)/C_G(H)\cong\operatorname{Im}\Phi\leqslant\operatorname{Aut}(H)$。
\end{proof}

$N/C$ 定理在有限群论中极为有用：它把正规化子/中心化子的商群嵌入 $\operatorname{Aut}(H)$，给出了群结构的强约束。例如，若 $H$ 的自同构群很小（如 $H$ 是循环群且 $\varphi(|H|)$ 小），则 $N_G(H)/C_G(H)$ 也很小，这限制了 $H$ 在 $G$ 中的"非中心"部分。

\medskip
\noindent\emph{$N/C$ 定理应用于 Sylow 子群。}
设 $P$ 是 $G$ 的 Sylow $p$-子群（定义见本附录首段及\cref{thm:sylow}）。由 Sylow 定理，$N_G(P)$ 控制 $P$ 的共轭类。$N/C$ 定理给出 $N_G(P)/C_G(P)\hookrightarrow\operatorname{Aut}(P)$。若 $P\cong C_p$（$p$ 阶循环群），则 $|\operatorname{Aut}(P)|=p-1$，故 $|N_G(P)/C_G(P)|\bigm|(p-1)$。又 $P\leqslant C_G(P)$（$P$ 是 Abel 群时 $P$ 中心化自身），故 $|N_G(P)/P|\bigm|(p-1)$。这限制了 Sylow $p$-子群的正规化子的大小。

特别地，若 $p$ 是 $|G|$ 的最高次幂且 $p\nmid(p-1)$（总是如此），则 $N_G(P)/C_G(P)$ 的阶既整除 $p-1$ 又是 $p$ 的幂（因 $N_G(P)/P$ 的阶整除 $|G|/|P|$ 与 $p$ 互质），故 $N_G(P)=C_G(P)$，即 Sylow $p$-子群在自身正规化子中是中心的——这就是 Burnside 正规 $p$-补定理的一个关键推论。
\par\medskip

\begin{theorem}[Frattini 论证]
\label{thm:frattini-argument}
设 $G$ 是有限群，$N\mathrel{\trianglelefteq} G$，$P$ 是 $N$ 的 Sylow $p$-子群。则
\[
G=N\cdot N_G(P)=\{ng\mid n\in N,\,g\in N_G(P)\}.
\]
\end{theorem}

\begin{proof}
对任意 $g\in G$，$gPg^{-1}\leqslant gNg^{-1}=N$（因 $N\mathrel{\trianglelefteq} G$），且 $|gPg^{-1}|=|P|$，故 $gPg^{-1}$ 也是 $N$ 的 Sylow $p$-子群。由 Sylow 定理（在 $N$ 中应用），存在 $n\in N$ 使 $n(gPg^{-1})n^{-1}=P$，即 $(ng)P(ng)^{-1}=P$，故 $ng\in N_G(P)$，即 $g=n^{-1}\cdot ng\in N\cdot N_G(P)$。
\end{proof}

Frattini 论证看似简单，但它是有限群论中最常用的技巧之一。它的典型用法是：若 $N\mathrel{\trianglelefteq} G$ 且 $N$ 的 Sylow 子群 $P$ 的正规化子 $N_G(P)$ 已知（例如 $P$ 是 $N$ 的特征子群，则 $N_G(P)\geqslant N$），则可以分解 $G$ 并推导结构信息。

\medskip
\noindent\emph{Frattini 论证的应用：$p$-群的作用。}
设 $P$ 是 $G$ 的正规 Sylow $p$-子群（即 $P\mathrel{\trianglelefteq} G$）。在 Frattini 论证中取 $N=P$，$P$ 自身是 $P$ 的唯一 Sylow $p$-子群，故 $N_G(P)=G$。论证给出 $G=P\cdot G$（平凡），但更有趣的是取 $N=P$ 的某个正规子群 $K\mathrel{\trianglelefteq} G$，$K\leqslant P$：$G=K\cdot N_G(Q)$，其中 $Q$ 是 $K$ 的 Sylow $p$-子群。若 $K$ 是 $p$-群，$Q=K$，$N_G(Q)\geqslant N_G(K)=G$，故 $G=K\cdot G$（仍平凡）。真正的力量体现在 $K$ 不是 $p$-群时：Frattini 论证把 $G$ 分解为 $K$ 和一个可能更小的群 $N_G(Q)$，可以用于归纳证明。
\par\medskip

\begin{proposition}[中心化子的正规化子性质]
\label{prop:centralizer-normalizer}
设 $H\leqslant G$。则：
\begin{enumerate}
\item $C_G(H)\mathrel{\trianglelefteq} N_G(H)$。
\item 若 $H$ 是 Abel 群，则 $H\leqslant C_G(H)$，且 $H\mathrel{\trianglelefteq} N_G(H)$。
\item 若 $H\mathrel{\trianglelefteq} G$ 且 $H$ 是 Abel 群，则 $H\leqslant Z(G)$ 当且仅当 $C_G(H)=G$。
\end{enumerate}
\end{proposition}

\begin{proof}
(1) 对 $n\in N_G(H)$ 和 $c\in C_G(H)$，需证 $ncn^{-1}\in C_G(H)$。对任意 $h\in H$：
\[
ncn^{-1}h(ncn^{-1})^{-1}=nc(n^{-1}hn)c^{-1}n^{-1}.
\]
因 $n\in N_G(H)$，$n^{-1}hn\in H$，而 $c\in C_G(H)$ 与 $H$ 的所有元素交换，故 $c(n^{-1}hn)c^{-1}=n^{-1}hn$。于是上式 $=n(n^{-1}hn)n^{-1}=h$。

(2) $H$ Abel 蕴含 $H$ 的元素彼此交换，故 $H\leqslant C_G(H)$。对 $n\in N_G(H)$，$nHn^{-1}=H$，这正是 $N_G(H)$ 的定义。

(3) $H\leqslant Z(G)$ 当且仅当 $G$ 的每个元素与 $H$ 的每个元素交换，即 $G\leqslant C_G(H)$，即 $C_G(H)=G$。
\end{proof}

\begin{theorem}[对称群的自同构群]
\label{thm:symmetric-aut}
对 $n\geqslant7$，$n\neq6$，
\[
\operatorname{Aut}(S_n)=\operatorname{Inn}(S_n)\cong S_n,
\]
即 $S_n$ 的所有自同构都是内自同构（$\operatorname{Out}(S_n)=1$）。唯一的例外是 $n=6$：$\operatorname{Out}(S_6)\cong C_2$，$S_6$ 有外自同构。
\end{theorem}

这个定理的证明较长（关键是证明每个自同构把对换映到对换，然后利用对换生成 $S_n$），但结论极为重要：$S_n$（$n\neq6$）的对称性"全部来自自身"——没有外部对称。$n=6$ 的外自同构是有限群论中最著名的例外之一，它与 $S_6$ 的特殊结构有关（$S_6$ 有 12 个 Sylow $5$-子群而非"通常"的 6 个，导致额外的对称性）。

\medskip
\noindent\emph{$S_6$ 的外自同构。}
$S_6$ 的外自同构 $\sigma\in\operatorname{Aut}(S_6)\setminus\operatorname{Inn}(S_6)$ 把对换（2-循环）映到\emph{三对对换的乘积}（即 $(ab)\mapsto(a'b')(c'd')(e'f')$，将 6 个字母分成三对）。这种映射不来自 $S_6$ 的任何共轭，但保持了 $S_6$ 的乘法结构。$S_6$ 的外自同构与数学中许多"偶然例外"现象相关，如 Cremona 变换、正二十面体的特殊对称等。
\par\medskip

\section{Jordan--H\"older 定理与可解群性质的完整证明}
\label{sec:jordan-holder-proof}

本节给出正文\cref{sec:composition-solvable} 中陈述的 Jordan--H\"older 定理与可解群基本性质的完整证明。

\par\medskip
\noindent\textbf{Jordan--H\"older 定理（\cref{thm:jordan-holder}）的证明}\quad
对 $|G|$ 归纳。若两条合成列的首个真子群相同，即 $G_1=H_1$，则对 $G_1$ 使用归纳假设即可。

设 $G_1\ne H_1$，并令 $N=G_1\cap H_1$。由于 $G/G_1$ 与 $G/H_1$ 都是单群，$G_1,H_1$ 是 $G$ 的极大正规子群。因此 $G_1H_1$ 只能是 $G$，而第二同构定理给出交叉同构
\[
G_1/N\cong G/H_1,\qquad H_1/N\cong G/G_1.
\]
取 $N$ 的任一合成列 $N=N_0\rhd\cdots\rhd N_k=\{e\}$。因为上式右端为单群，
\[
G_1\rhd N\rhd N_1\rhd\cdots\rhd\{e\},\qquad
H_1\rhd N\rhd N_1\rhd\cdots\rhd\{e\}
\]
分别是 $G_1$ 与 $H_1$ 的合成列。对这两个较小的群应用归纳假设，可用它们替换原合成列的尾部。替换后，两边共有 $N$ 的全部合成因子；其余两项分别为
\[
\{G/G_1,\,G_1/N\},\qquad \{G/H_1,\,H_1/N\},
\]
而交叉同构表明这两个二元多重集相同。因此两条原合成列的合成因子仅相差排列，长度也相同。
\par\medskip

\begin{proposition}[可解群的基本性质]
\begin{enumerate}[label=(\roman*),leftmargin=3em]
\item 可解群的子群和商群可解。
\item 可解群的可解扩张可解：若 $N\mathrel{\trianglelefteq} G$，$N$ 和 $G/N$ 都可解，则 $G$ 可解。
\item $p$-群可解（因 $p$-群有中心链，每步商为 $C_p$）。
\item $S_n$ 对 $n\ge5$ 不可解（因 $A_n$ 是单群且非 Abel）。
\end{enumerate}
\end{proposition}

\par\medskip
\noindent\textbf{可解群性质的证明}\quad
(i) 子群可解. 设 $H\le G$，$G$ 有 Abel 链 $G=G_0\rhd G_1\rhd\cdots\rhd G_n=\{e\}$。令 $H_i=H\cap G_i$。则
\[
H=H_0\ge H_1\ge\cdots\ge H_n=\{e\}.
\]
需验证每步正规且商 Abel。$H_{i+1}=H\cap G_{i+1}\mathrel{\trianglelefteq} H\cap G_i=H_i$：对 $h\in H_i$，$hG_{i+1}h^{-1}=G_{i+1}$（因 $G_{i+1}\mathrel{\trianglelefteq} G_i$ 且 $h\in G_i$），故 $hH_{i+1}h^{-1}=h(H\cap G_{i+1})h^{-1}=H\cap hG_{i+1}h^{-1}=H\cap G_{i+1}=H_{i+1}$。

商群 $H_i/H_{i+1}$ 嵌入 Abel 群 $G_i/G_{i+1}$：定义
\[
\psi:H_i\to G_i/G_{i+1},\qquad\psi(h)=hG_{i+1}.
\]
$\operatorname{Ker}\psi=H_i\cap G_{i+1}=H\cap G_i\cap G_{i+1}=H\cap G_{i+1}=H_{i+1}$。由第一同构定理 $H_i/H_{i+1}\hookrightarrow G_i/G_{i+1}$，Abel 群的子群 Abel，故 $H$ 可解。

(ii) 商群可解. 设 $N\mathrel{\trianglelefteq} G$，考虑投影 $\pi:G\to G/N$，令 $Q_i=\pi(G_i)=G_iN/N$。则
\[
G/N=Q_0\ge Q_1\ge\cdots\ge Q_n=\{eN\}.
\]
$Q_{i+1}\mathrel{\trianglelefteq} Q_i$ 因 $G_{i+1}\mathrel{\trianglelefteq} G_i$。商 $Q_i/Q_{i+1}$ 是 $G_i/G_{i+1}$ 的商群（由对应定理），Abel 群的商群 Abel。故 $G/N$ 可解。

(iii) 扩张可解. 设 $N\mathrel{\trianglelefteq} G$，$N$ 有 Abel 链 $N=N_0\rhd N_1\rhd\cdots\rhd N_k=\{e\}$，$G/N$ 有 Abel 链 $G/N=Q_0\rhd Q_1\rhd\cdots\rhd Q_m=\{eN\}$。令 $G_i=\pi^{-1}(Q_i)$（$\pi:G\to G/N$）。由对应定理 $G_i\mathrel{\trianglelefteq} G_{i-1}$ 且 $G_{i-1}/G_i\cong Q_{i-1}/Q_i$ 是 Abel。拼接两条链：
\[
G=G_0\rhd G_1\rhd\cdots\rhd G_m=N\rhd N_1\rhd\cdots\rhd N_k=\{e\},
\]
每个商 Abel，故 $G$ 可解。

(iv) $p$-群可解. 设 $|G|=p^a$。由类方程与\cref{cor:pgroup-center}，中心 $Z(G)$ 非平凡。取 $Z(G)$ 中 $p$ 阶子群 $C_p\le Z(G)$（Cauchy 定理）。因 $C_p\le Z(G)$，$C_p\mathrel{\trianglelefteq} G$。商群 $G/C_p$ 的阶为 $p^{a-1}$，仍是 $p$-群。由归纳假设 $G/C_p$ 可解，$C_p$ 亦可解，再由 (iii) 知 $G$ 可解。归纳基础为 $|G|=p^0=1$ 时 $G=\{e\}$ 可解（空链）。

(v) $S_n$ 不可解（$n\ge5$）. $A_n$（$n\ge5$）是单群且非 Abel（\cref{ex:a5-simplicity} 对 $n=5$ 的证明可推广到一般 $n\ge5$）。若 $S_n$ 可解，则由 (i) 的子群遗传性，$A_n\le S_n$ 也应可解。但 $A_n$ 的合成列只能是
\[
A_n\rhd\{e\},
\]
因 $A_n$ 单。合成因子 $A_n$ 非 Abel，不满足可解群定义。矛盾，故 $S_n$ 不可解。
\par\medskip

\section{幂零群}
\label{sec:nilpotent}

可解群要求合成因子是 Abel 群，但不要求这些 Abel 商如何嵌入原群。幂零群进一步加强：要求每一步商不仅 Abel，而且以\emph{中心}方式嵌入。这使得幂零群的结构比一般可解群更"刚性"。

\begin{definition}[幂零群与中心列]
群 $G$ 是\emph{幂零群}，若存在子群链
\[
G=G_0\ge G_1\ge\cdots\ge G_n=\{e\},
\]
使得 $G_{i-1}/G_i\le Z(G/G_i)$ 对所有 $i$（即 $[G,G_{i-1}]\le G_i$）。这样的链称为\textbf{中心列}，最短长度称为\textbf{幂零类} $c(G)$。
\end{definition}

条件 $G_{i-1}/G_i\le Z(G/G_i)$ 比 $G_{i-1}/G_i$ Abel 更强：它要求 $G_{i-1}/G_i$ 不仅交换，而且与整个 $G/G_i$ 交换。等价地，$[G,G_{i-1}]\le G_i$：换位子 $[g,h]$（$g\in G$，$h\in G_{i-1}$）落入 $G_i$。

\begin{definition}[上中心列与下中心列]
$G$ 的\textbf{上中心列}定义为
\[
\{e\}=Z_0\le Z_1\le Z_2\le\cdots,\qquad Z_{i+1}/Z_i=Z(G/Z_i),
\]
即 $Z_1=Z(G)$，$Z_{i+1}$ 是 $G/Z_i$ 的中心在 $G$ 中的原像。\textbf{下中心列}定义为
\[
G=\gamma_1\ge\gamma_2\ge\cdots,\qquad\gamma_{i+1}=[G,\gamma_i],
\]
即 $\gamma_2=[G,G]=G'$（导群），$\gamma_3=[G,G']$，等等。
\end{definition}

\begin{theorem}[幂零群的等价刻画]
\label{thm:nilpotent-equivalent}
对有限群 $G$，以下等价：
\begin{enumerate}[label=(\roman*),leftmargin=3em]
\item $G$ 幂零。
\item 上中心列终止于 $G$：存在 $c$ 使 $Z_c(G)=G$。
\item 下中心列终止于 $\{e\}$：存在 $c$ 使 $\gamma_{c+1}(G)=\{e\}$。
\item $G$ 是其 Sylow 子群的直积：$G=P_1\times P_2\times\cdots\times P_k$。
\end{enumerate}
满足上述条件时，幂零类 $c(G)$ 等于使 $Z_c=G$ 的最小 $c$，也等于使 $\gamma_{c+1}=\{e\}$ 的最小 $c$。
\end{theorem}

\begin{proof}
(i)$\Leftrightarrow$(ii). 若 $G$ 有中心列 $G=G_0\ge\cdots\ge G_n=\{e\}$，则 $G_i\le Z_{n-i}$（对 $n-i$ 归纳：$G_n=\{e\}=Z_0$；若 $G_i\le Z_{n-i}$，则 $[G,G_{i-1}]\le G_i\le Z_{n-i}$，故 $G_{i-1}/Z_{n-i}\le Z(G/Z_{n-i})$，即 $G_{i-1}\le Z_{n-i+1}$）。故 $G=G_0\le Z_n$，$Z_n=G$。逆方向类似。

(ii)$\Leftrightarrow$(iii). 对 $i$ 归纳可证 $\gamma_{i+1}\le Z_{c-i}$（$c$ 为上中心列长度）：$\gamma_1=G=Z_c$；若 $\gamma_i\le Z_{c-i+1}$，则 $\gamma_{i+1}=[G,\gamma_i]\le[G,Z_{c-i+1}]\le Z_{c-i}$（因 $Z_{c-i+1}/Z_{c-i}=Z(G/Z_{c-i})$ 蕴含 $[G,Z_{c-i+1}]\le Z_{c-i}$）。故 $\gamma_{c+1}\le Z_0=\{e\}$。逆方向类似。

(iii)$\Leftrightarrow$(iv). 这是有限群论的经典结果。\textbf{(iv)$\Rightarrow$(iii):} 直积的换位子 $[P_1\times\cdots\times P_k,P_1\times\cdots\times P_k]=[P_1,P_1]\times\cdots\times[P_k,P_k]$（不同 Sylow 分量交换），每个 $p$-群有下中心列终止于 $\{e\}$（因 $p$-群的中心非平凡，上中心列严格增长直到全群），故直积也如此。\textbf{(i)$\Rightarrow$(iv):} 幂零群的每个 Sylow 子群正规（因 Sylow $p$-子群在幂零群中是唯一的：若 $P,Q$ 都是 Sylow $p$-子群，$P\le Z_c$ 的某层但不在 $Z_{c-1}$，则 $Q$ 正规化 $P$，由 Sylow 共轭性 $P=Q$）。不同 Sylow 子群的交为 $\{e\}$（阶互质），元素彼此交换，故 $G$ 是直积。
\end{proof}

\begin{proposition}[幂零群的基本性质]
\label{prop:nilpotent-properties}
\begin{enumerate}[label=(\roman*),leftmargin=3em]
\item 幂零群的子群和商群幂零。
\item 幂零群的直积幂零，且 $c(G\times H)=\max(c(G),c(H))$。
\item $p$-群幂零，幂零类 $c\le\log_p|G|-1$。
\item 幂零群可解，但可解群不一定幂零。
\item 幂零群的每个 Sylow 子群正规（唯一 Sylow $p$-子群）。
\item 有限幂零群中，正规化子条件：$H\lneq N_G(H)$ 对所有真子群 $H$ 成立。
\end{enumerate}
\end{proposition}

\begin{proof}
(i) 子群：中心列与子群的交给出子群的中心列。商群：中心列取商给出商的中心列。

(ii) 直积的上中心列 $(Z_i(G\times H))=Z_i(G)\times Z_i(H)$，故 $c(G\times H)=\max(c(G),c(H))$。

(iii) $p$-群有非平凡中心（\cref{cor:pgroup-center}），$Z_1=Z(G)\ne\{e\}$，$G/Z_1$ 仍是 $p$-群故有非平凡中心，$Z_2/Z_1=Z(G/Z_1)\ne\{e\}$，$Z_2$ 严格大于 $Z_1$。每步 $|Z_{i+1}/Z_i|\ge p$，故 $c\le\log_p|G|-1$。

(iv) 中心列的商 $G_{i-1}/G_i\le Z(G/G_i)$ 是 Abel 的，故中心列也是 Abel 链，幂零蕴含可解。反例：$S_3$ 可解（Abel 链 $S_3\rhd A_3\rhd\{e\}$）但不幂零（$Z(S_3)=\{e\}$，上中心列停滞在 $\{e\}$；或等价地，$S_3$ 的 Sylow $2$-子群不正规）。

(v) 由 \cref{thm:nilpotent-equivalent} (iv)，$G$ 是 Sylow 子群直积，每个因子唯一故正规。

(vi) 设 $H\lneq G$，取中心列 $G=Z_c\ge\cdots\ge Z_0=\{e\}$。设 $i$ 最大使 $Z_i\le H$（存在因 $Z_0=\{e\}\le H$，$Z_c=G\nleq H$）。则 $Z_{i+1}\nleq H$，但 $[G,Z_{i+1}]\le Z_i\le H$，故对 $z\in Z_{i+1}\setminus H$ 和 $h\in H$，$[z,h]=zhz^{-1}h^{-1}\in Z_i\le H$，即 $zhz^{-1}\in H$，$z\in N_G(H)$。故 $H\lneq\langle H,z\rangle\le N_G(H)$。
\end{proof}

\medskip
\noindent\emph{幂零群的实例与非实例。}
幂零:
\begin{itemize}
\item 所有 $p$-群：如 $D_4$（$2^3$ 阶），上中心列 $\{e\}\le\langle r^2\rangle\le\langle r\rangle\le D_4$，幂零类 $2$。
\item $D_4$ 幂零但 $D_3\cong S_3$ 不幂零：$D_3$ 的 Sylow $2$-子群（三个反射子群）不正规。
\item Abel 群幂零类 $0$ 或 $1$（$Z_1=G$）。
\item 直积 $C_4\times C_3\cong C_{12}$ 幂零（Sylow 分解 $C_4\times C_3$）。
\end{itemize}
不幂零:
\begin{itemize}
\item $S_3$：可解但不幂零（Sylow $2$-子群不正规）。
\item $S_4$：可解但不幂零（Sylow $3$-子群不正规）。
\item $A_4$：可解但不幂零（Sylow $2$-子群 $V_4$ 正规，但 Sylow $3$-子群不正规）。
\end{itemize}
\par\medskip

\medskip
\noindent\emph{Heisenberg 群：幂零类 $2$ 的典型例子。}
上三角整数矩阵群
\[
H_3(\ZZ)=\left\{\begin{pmatrix}1&a&c\\0&1&b\\0&0&1\end{pmatrix}:a,b,c\in\ZZ\right\}
\]
是幂零类 $2$ 的幂零群。先把换位子真正算出来。记
\[
x=\begin{pmatrix}1&a&c\\0&1&b\\0&0&1\end{pmatrix},
\qquad
y=\begin{pmatrix}1&a'&c'\\0&1&b'\\0&0&1\end{pmatrix},
\]
直接相乘得
\[
xy=\begin{pmatrix}1&a+a'&c+c'+ab'\\0&1&b+b'\\0&0&1\end{pmatrix},
\qquad
yx=\begin{pmatrix}1&a+a'&c+c'+a'b\\0&1&b+b'\\0&0&1\end{pmatrix},
\]
两式仅右上角的 $c$ 分量相差 $ab'-a'b$；于是
\[
[x,y]=xyx^{-1}y^{-1}
=\begin{pmatrix}1&0&ab'-a'b\\0&1&0\\0&0&1\end{pmatrix},
\]
换位子只在 $c$ 分量非零。由此
\[
Z(H_3)=\left\{\begin{pmatrix}1&0&c\\0&1&0\\0&0&1\end{pmatrix}:c\in\ZZ\right\}\cong\ZZ,
\]
导群 $[H_3,H_3]=Z(H_3)$，故下中心列 $H_3\ge[H_3,H_3]=Z(H_3)\ge[H_3,Z(H_3)]=\{I\}$ 终止于第二步。上中心列 $\{I\}\le Z(H_3)\le H_3$ 也两步终止。

物理中，Heisenberg 群是量子力学位置--动量对易关系 $[x,p]=\ii\hbar$ 的群论化身：位置和动量的平移生成 $H_3$，中心分量是相空间面积（几何相位）。幂零性反映 $[x,p]$ 是中心元素（与一切交换），故"二次换位子"自动消失。
\par\medskip

\section{自由群与 Tietze 变换}
\label{sec:free-group-tietze}

正文中"生成元与关系"的写法 $G=\langle S\mid R\rangle$ 依赖自由群的概念；这里给出正式定义。

\begin{definition}[自由群]
集合 $S=\{s_1,\dots,s_n\}$ 上的自由群 $F(S)$ 是所有形式为
\[
s_{i_1}^{\epsilon_1}s_{i_2}^{\epsilon_2}\cdots s_{i_k}^{\epsilon_k},
\qquad\epsilon_j=\pm1,
\]
的\emph{既约字}（无相邻 $s_i s_i^{-1}$ 或 $s_i^{-1}s_i$ 可消去）构成的群，乘法为拼接后消去。$F(S)$ 满足泛性质：对任意群 $G$ 和任意映射 $\varphi:S\to G$，存在唯一同态 $\bar\varphi:F(S)\to G$ 扩展 $\varphi$。
\end{definition}

因此展示 $G=\langle S\mid R\rangle$ 的准确含义是 $G\cong F(S)/\langle\!\langle R\rangle\!\rangle$，其中 $\langle\!\langle R\rangle\!\rangle$ 是关系集 $R$ 生成的正规子群。

群的展示并不唯一。Tietze 变换给出两类保持群不变的改写：
\begin{itemize}
\item \emph{I 型}：添加一个可由现有关系导出的关系（或删除冗余关系）；
\item \emph{II 型}：添加一个新生成元 $a$ 及关系 $a=w$（$w$ 是旧生成元的字），或删除这样的生成元和关系。
\end{itemize}
正文把 $D_3=\langle r,s\mid r^3,s^2,srsr\rangle$ 改写为 Coxeter 形式 $\langle a,b\mid a^2,b^2,(ab)^3\rangle$ 正是 II 型变换（引入 $a=s$、$b=sr$）的实例；任意两个有限展示之间总存在有限的 Tietze 变换序列，这是展示理论的基本定理（证明从略）。

\section{\texorpdfstring{$A_5$}{A5} 单性的完整枚举}
\label{sec:a5-enumeration}

正文用一句话给出 $A_5$ 单性的计数理由；这里把逐一枚举补齐。
\label{ex:a5-simplicity}

交错群 $A_5$ 有 $|A_5|=60=2^2\cdot3\cdot5$。用 Sylow 定理可以证明 $A_5$ 是单群（没有非平凡正规子群），这是 Galois 理论中"五次方程一般不可解"的核心事实。

先由 Sylow 定理核对各素数幂阶子群的数目：
\begin{itemize}
\item Sylow $5$-子群：$n_5\equiv1\pmod5$ 且 $n_5\mid12$。$A_5$ 中全部 $24$ 个 $5$-循环都是偶置换，每个 $C_5$ 含 $4$ 个非单位元，故 $n_5=6$。
\item Sylow $3$-子群：$n_3\equiv1\pmod3$ 且 $n_3\mid20$，故 $n_3\in\{1,4,10\}$。$A_5$ 中有 $20$ 个 $3$-循环（$(abc)$ 型），每个 $C_3$ 含 $2$ 个，故 $n_3=10$。
\item Sylow $2$-子群：$n_2\equiv1\pmod2$ 且 $n_2\mid15$，故 $n_2\in\{1,3,5,15\}$。$A_5$ 中 Klein 四元群 $V_4=\{e,(12)(34),(13)(24),(14)(23)\}$ 是 Sylow $2$-子群（$4$ 阶），共 $5$ 个（对应 $5$ 种把 $\{1,2,3,4,5\}$ 分成一对加单点的方式），故 $n_2=5$。
\end{itemize}

另一方面，正规子群必须是共轭类的并。$A_5$ 的共轭类大小为
\[
1,\quad 12,\quad 12,\quad 15,\quad 20,
\]
分别对应单位元、$20$ 个 $3$-循环、分成两类的 $24$ 个 $5$-循环，以及 $15$ 个双对换。

若 $N\mathrel{\trianglelefteq} A_5$，$|N|\mid60$ 且 $|N|=1+a\cdot20+b\cdot12+c\cdot12+d\cdot15$（$a,b,c,d\in\{0,1\}$，因每类要么全在要么全不在）。逐一检查 $|N|\in\{1,2,3,4,5,6,10,12,15,20,30,60\}$：
\begin{itemize}
\item $|N|=1$：$N=\{e\}$。
\item $|N|=60$：$N=A_5$。
\item $|N|=5$：需 $1+d\cdot15=5$ 或 $1+c\cdot12=5$ 等，均无整数解。
\item $|N|=10$：$1+a\cdot20=10$ 不行；$1+c\cdot12=10$ 不行；$1+d\cdot15=10$ 不行；$1+b\cdot12+c\cdot12=10$ 不行。
\item $|N|=15$：$1+d\cdot15=16\neq15$；$1+a\cdot20=15$ 不行；$1+b\cdot12=15$ 不行。
\item $|N|=20$：$1+a\cdot20=21\neq20$；$1+d\cdot15=20$ 不行；$1+b\cdot12=20$ 不行。
\item $|N|=30$：$1+a\cdot20+d\cdot15=30$ 即 $20a+15d=29$，无解；$1+b\cdot12+c\cdot12=30$ 即 $12(b+c)=29$，无解；$1+a\cdot20=30$ 不行。
\item $|N|=12$：$1+b\cdot12=13\neq12$；$1+c\cdot12=13\neq12$。
\item $|N|=6$：$1+a\cdot20=6$ 不行；$1+d\cdot15=6$ 不行。
\item $|N|=4$：$1+d\cdot15=4$ 不行。
\item $|N|=3$：$1+a\cdot20=3$ 不行。
\item $|N|=2$：$1+d\cdot15=2$ 不行。
\end{itemize}
因此 $A_5$ 没有非平凡正规子群，是单群。

$A_5$ 的单性意味着 $S_5$ 的合成列为
\[
S_5\rhd A_5\rhd\{e\},
\]
合成因子 $C_2$ 和 $A_5$。因 $A_5$ 非交换单群，$S_5$ 不可解。由 Galois 理论，五次方程一般不可根式求解。物理上，$A_5\cong I$（正二十面体转动群）的单性意味着正二十面体没有非平凡的连续对称性约化——它的全部对称性"不可拆分"。
